\documentclass[11pt,a4paper]{article}

\usepackage[utf8]{inputenc}
\usepackage[T1]{fontenc}
\usepackage[norsk]{babel}
\usepackage[a4paper, margin=2.4cm]{geometry}
\usepackage{hyperref}
\usepackage{enumitem}
\usepackage{longtable}

\setlength{\parindent}{0pt}
\setlength{\parskip}{0.8em}

\title{Enkel brukerguide for \texttt{python/main.py}}
\author{TDT4145 prosjekt}
\date{19. mars 2026}

\begin{document}

\maketitle

\section*{Slik brukes programmet}
Denne guiden er laget for en ekstern tester som skal kjøre prosjektet via
\texttt{python/main.py}. Programmet er menybasert, og hvis et valg gir en
forventet feilmelding, kommer menyen opp igjen i stedet for at programmet
stopper.

Start i prosjektmappen og kjør \texttt{python3 python/main.py}. Velg først
\texttt{1. Legg inn data (SQL-script)}. Dette legger inn manglende demo-data og
kan kjøres flere ganger uten at hele databasen bygges på nytt.

Deretter kan programmet demonstreres i denne rekkefølgen:

\begin{enumerate}[leftmargin=1.6em]
    \item Velg \texttt{4. Ukeplan}. Skriv for eksempel \texttt{2026-03-19} som
    startdag og \texttt{12} som uke. Da får du oversikt over treningene med dato,
    tidspunkt og senter.
    \item Velg \texttt{2. Book trening} hvis du vil legge Johnny på en fremtidig
    time. Et eksempel per 19. mars 2026 er \texttt{johnny@stud.ntnu.no},
    aktivitet \texttt{HIIT}, dato \texttt{2026-03-20}, starttid \texttt{07:00},
    senter \texttt{Moholt}.
    \item Velg \texttt{3. Registrer oppmøte}. Her skriver du bare epost, og
    programmet viser hvilke treninger brukeren er påmeldt. Du trenger altså ikke
    å vite trenings-ID; du velger via treningsnavn fra listen.
    \item Velg \texttt{5. Personlig besøkshistorie} med
    \texttt{johnny@stud.ntnu.no} og \texttt{2026-01-01}.
    \item Velg \texttt{6. Svartelisting} med \texttt{mikkels@stud.ntnu.no}.
    \item Velg \texttt{7. Mest aktive bruker} med måned \texttt{3}.
    \item Velg \texttt{8. Finn treningspartnere}. Dette valget trenger ingen
    ekstra input.
\end{enumerate}

\section*{Viktig å vite}
I brukstilfelle 6 er det med vilje valgt å bruke \texttt{mikkels@stud.ntnu.no}
og ikke Johnny. Grunnen er at Johnny brukes i flere av de andre brukstilfellene,
og det er derfor mer praktisk å holde ham tilgjengelig der. Dette er altså et
bevisst valg i demonstrasjonen.

Hvis du prøver å registrere oppmøte for Johnny på en time som allerede er ferdig,
for eksempel en time 17. mars når dagens dato er 19. mars 2026, skal det feile.
Det er meningen. Systemet skal ikke registrere oppmøte bakover i tid. Hvis du vil
se et vellykket oppmøte, finner du først en fremtidig gruppetime i
\texttt{4. Ukeplan}, booker den i \texttt{2. Book trening}, og registrerer
deretter oppmøte via \texttt{3. Registrer oppmøte}.

\end{document}
